\documentclass[11pt]{article}
\usepackage{mathblatt}
\begin{document}
\blattkopf*{Basisaufgaben}{BAS-Z9}{Basisaufgaben · Zettel · BAS-Z9}
\einheitenkopf[][Zettel 9]{Basisaufgaben · Zettel 9}
\zweigzeile{ohne Taschenrechner · je Aufgabe 1 Punkt · Lösungen auf der Rückseite}
% \small: zehn Aufgaben mit bis zu zwei Koordinatensystemen auf einer Seite (Render 28.09.)
\small

% 1: Bruchteil einer Fläche bestimmen – brueche-dezimalzahlen-basis-k1-v8 (Original 2026-FOR-B1b)
\begin{aufgabe}{Bei welcher Figur ist genau $\frac{5}{6}$ grau? \hfill (P26F) \\ \kreuz{P} \kreuz{Q} \kreuz{R} \kreuz{S}}

P \kreissektor[0.8]{300}{} \quad Q \bruchrechteck[2.5]{5}{8} \quad R \bruchkreis[0.8]{4}{5} \quad S \bruchrechteck[2.5]{4}{6}
\end{aufgabe}

% 2: Bruchteil einer Größe berechnen – brueche-dezimalzahlen-basis-k2-v3 (Original 2018-OS-B1a)
\begin{aufgabe}{$\frac{2}{5}$ von $3$ l – wie viel Liter? \hfill (P18) \\ \leerfeld[l]}
\end{aufgabe}

% 3: Zahl zu Bedingung angeben – rationale-zahlen-basis-k3-v2 (Original 2015-OS-B1b)
\begin{aufgabe}{Gib eine Zahl an, die kleiner ist als $-30$. \hfill (P15) \\ \leerfeld}
\end{aufgabe}

% 4: Wurzel eines Quadrats berechnen – potenzen-wurzeln-basis-k2-v1 (Original 2015-OS-B1j)
\begin{aufgabe}{Welche Aussage ist wahr? Kreuze an. \hfill (P15*)\\ \kreuz{$\sqrt{(-5)^2} = -5$}\\ \kreuz{$\sqrt{(-5)^2} = 5$}\\ \kreuz{$\sqrt{(-5)^2}$ ist nicht definiert}}
\end{aufgabe}

% 5: Zehnerpotenzschreibweise umwandeln – potenzen-wurzeln-basis-k3-v4 (Original 2025-OS-B1d)
\begin{aufgabe}{$510 = 5{,}1 \cdot 10^{\square}$ – welche Hochzahl? \hfill (P25) \\ \leerfeld}
\end{aufgabe}

% 6: Prozentwert berechnen – prozentrechnung-basis-k4-v6 (Original 2026-FOR-B1a)
\begin{aufgabe}{$15\,\%$ von $80\,€$ \hfill (P26F) \\ \leerfeld[€]}
\end{aufgabe}

% 7: Antiproportionale Zuordnung Dreisatz – zuordnungen-basis-k1-v2 (Original 2014-GYM-B1c)
\begin{aufgabe}{3 Maler streichen ein Haus in 8 Tagen. Wie viele Tage brauchen 4 Maler? \hfill (P14G) \\ \leerfeld[Tage]}
\end{aufgabe}

% 8: Lineare Gleichung lösen – lineare-gleichungen-basis-k1-v5 (Original 2024-OS-B1d)
\begin{aufgabe}{Löse die Gleichung $6 \cdot (x - 1{,}5) = 0$. \hfill (P24) \\ x = \leerfeld}
\end{aufgabe}

% 9: Wertetabelle einer Funktion zuordnen – quadratische-funktionen-basis-k1-v1 (Original 2026-FOR-B1e)
\begin{aufgabe}{Welche Wertetabelle gehört zu $y = 2x^2$? \hfill (P26F) \\ \kreuz{A} \kreuz{B} \kreuz{C}}

A: \wertetabelle[-4,-2,0,2,4]{x}{y}{-2,-1,0,1,2} B: \wertetabelle[16,4,0,4,16]{x}{y}{-2,-1,0,1,2} C: \wertetabelle[8,2,0,2,8]{x}{y}{-2,-1,0,1,2}
\end{aufgabe}

% 10: Winkelfunktion Seitenverhältnis angeben – trigonometrie-basis-k2-v5 (Original 2025-OS-B1g)
\begin{aufgabe}{\begin{minipage}[t]{0.6\linewidth}Rechter Winkel bei $A$. Trage den Bruch ein. \hfill (P25) \\ $\mathrm{tan}\,\beta =$ \leerfeld\end{minipage}\hfill\begin{minipage}[t]{0.37\linewidth}\vspace{0pt}
\dreieck{(0,0)}{(3,0)}{(0,2.2)}{d}{e}{f}{}{\beta}{}
\end{minipage}}
\end{aufgabe}

\begleitteil
\erg{1}{P, denn $\frac{300}{360} = \frac{5}{6}$}
\erg{2}{$3 : 5 \cdot 2 = 1{,}2$ l}
\erg{3}{z. B. $-40$; richtig ist jede Zahl links von $-30$, nicht $-20$}
\erg{4}{$\sqrt{(-5)^2} = 5$, denn erst das Quadrat: $(-5)^2 = 25$, und die Wurzel ist nie negativ}
\erg{5}{$2$ – das Komma wandert zwei Stellen nach links}
\erg{6}{$0{,}15 \cdot 80 = 12\,€$}
\erg{7}{$3 \cdot 8 = 24$ Malertage; $24 : 4 = 6$ Tage}
\erg{8}{$x - 1{,}5 = 0$, also $x = 1{,}5$}
\erg{9}{Tabelle C: $2 \cdot 2^2 = 8$ – erst quadrieren, dann malnehmen}
\erg{10}{$\mathrm{tan}\,\beta = \frac{e}{f}$}
\end{document}
